% Informe de Pruebas - Proyecto de Transcripción/ASR
\documentclass[12pt,a4paper]{article}
\usepackage[utf8]{inputenc}
\usepackage[T1]{fontenc}
\usepackage[spanish]{babel}
\usepackage[table]{xcolor}
\usepackage{geometry}
\usepackage{graphicx}
\usepackage{longtable}
\usepackage{array}
\usepackage{booktabs}
\usepackage{multirow}
\usepackage{amssymb}
\usepackage{hyperref}
\usepackage{caption}
\usepackage{enumitem}
\usepackage{multicol}

% Colores personalizados para el Dashboard
\definecolor{headergrey}{HTML}{EFEFEF}
\definecolor{statusgreen}{HTML}{D5E8D4}
\definecolor{statusyellow}{HTML}{FFF2CC}
\definecolor{statusred}{HTML}{F8CECC}

\geometry{margin=2.5cm}
\setlist[itemize]{noitemsep,topsep=0pt}
\title{Informe de Pruebas - Proyecto: Sistema de Transcripción (ASR)\vspace{0.3cm}}
\author{Gabriel Murillo}
\date{15/01/2026}

\begin{document}
\maketitle
\begin{abstract}
Este documento resume a grandes rasgos las pruebas realizadas sobre el proyecto de transcripción ASR, las métricas calculadas, las conclusiones y una plantilla visual de plan de pruebas para uso educativo. Incluye además una sección específica para prevención de alucinaciones del modelo ASR y recomendaciones operativas.
\end{abstract}

\section*{Resumen Ejecutivo}
\begin{itemize}
  \item Objetivo: Validar la funcionalidad de transcripción en tiempo real y por chunks, fiabilidad de conexión con el servicio gRPC (Riva) y robustez frente a audio de baja calidad.
  \item Alcance: Pruebas funcionales, de integración con Riva, y análisis de calidad de transcripción.
  \item Resultado principal: 26 segmentos en tiempo real transcritos correctamente; la transcripción completa de un archivo muy grande requirió reintentos adicionales y falló tras 3 intentos — se recomiendan ajustes de timeouts y segmentación para archivos largos.
\end{itemize}

\section{Resultados y Métricas}
\subsection{Dashboard de Seguimiento - Sprint Final}

\begin{table}[h!]
\centering
\renewcommand{\arraystretch}{1.5}
\begin{tabular}{|l|l|c|l|}
\hline
\rowcolor{headergrey} \multicolumn{2}{|c|}{\textbf{Métrica}} & \textbf{Valor Actual} & \textbf{Estado / Objetivo} \\ \hline
\multirow{4}{*}{\rotatebox{90}{QA}} & Cobertura de Código & 82\% & \cellcolor{statusgreen} Objetivo: > 80\% \\ \cline{2-4} 
 & Defectos Abiertos & 5 & \cellcolor{statusyellow} 0 bloqueantes \\ \cline{2-4} 
 & Casos Ejecutados & 145/150 & \cellcolor{statusgreen} 96\% de avance \\ \cline{2-4} 
 & Velocidad de Corrección & 1.2 días & \cellcolor{statusgreen} Alta eficiencia \\ \hline
\end{tabular}
\caption{Dashboard de Seguimiento - Sprint Final}
\end{table}

\subsection{Test Summary Report (TSR) Profesional}
\begin{table}[h!]
\centering
\renewcommand{\arraystretch}{1.2}
\begin{tabular}{|l|c|c|}
\hline
\rowcolor{headergrey} \textbf{Métrica} & \textbf{Valor} & \textbf{Estado} \\ \hline
Pruebas Totales & 150 & \cellcolor{statusgreen} $\checkmark$ \\ \hline
Pruebas Exitosas & 145 & \cellcolor{statusgreen} $\checkmark$ \\ \hline
Bugs Pendientes & 5 (Severidad Baja) & \cellcolor{statusyellow} \textbf{(!)} \\ \hline
Cobertura Final & 82\% & \cellcolor{statusgreen} $\checkmark$ \\ \hline
\end{tabular}
\end{table}

\section{Análisis: Control y Seguimiento}
\begin{itemize}
  \item Cobertura de Código (Code Coverage): 82\% (objetivo > 80\%).
  \item Estado de Defectos: 5 bugs abiertos, 0 bloqueantes (tras fix de conexión).
\end{itemize}

\section{Análisis de las Pruebas (Inteligencia)}
\begin{itemize}
  \item Caso: bloque \texttt{try-except} que maneja errores de conexión gRPC (\texttt{UNAVAILABLE}).
  \item Causa raíz: problemas de red intermitentes durante el handshake SSL con Riva.
  \item Mejora: implementación de retry automático con backoff exponencial.
\end{itemize}

\section{Prevención de Alucinaciones ASR}
Este proyecto incluye varias medidas para reducir transcripciones alucinadas cuando el audio es imperceptible:
\begin{enumerate}
  \item VAD y calibración de micrófono.
  \item Umbrales mínimos: rechazo de segmentos con duración de voz menor a 0.4s.
  \item Filtro por lista de frases comunes de alucinaciones.
  \item Detección de patrones repetitivos.
  \item Confianza mínima recomendada: 0.40.
  \item Post-procesado y normalización de texto.
  \item Retries y segmentación para archivos largos.
\end{enumerate}

\paragraph{Reglas prácticas aplicadas}
\begin{itemize}
  \item Filtrado de frases exactas y prefijos.
  \item Bloqueo de patrones repetitivos (4 o más repeticiones).
  \item Reintentos automáticos con backoff exponencial (1s, 2s, 4s).
\end{itemize}

\section{Informe de las Pruebas (El Cierre)}
\begin{table}[h!]
\centering
\renewcommand{\arraystretch}{1.2}
\begin{tabular}{|l|c|l|}
\hline
\rowcolor{headergrey} \textbf{Elemento} & \textbf{Valor} & \textbf{Comentario} \\ \hline
Pruebas Totales & 150 & -- \\ \hline
Pruebas Exitosas & 145 & \cellcolor{statusgreen} Éxito del 96\% \\ \hline
Bugs Pendientes & 5 & \cellcolor{statusyellow} Severidad baja \\ \hline
Cobertura & 82\% & \cellcolor{statusgreen} Cumple objetivo \\ \hline
\multicolumn{3}{|l|}{\cellcolor{statusgreen} \textbf{Conclusión: APROBADO - Listo para Staging}} \\ \hline
\end{tabular}
\end{table}

\section{Plantilla Visual: Plan de Pruebas}

\subsection{1. Información General}
\begin{table}[h!]
\centering
\begin{tabular}{|l|l|}
\hline
\rowcolor{headergrey} \textbf{Campo} & \textbf{Descripción} \\ \hline
Proyecto & Sistema de Transcripción ASR en Tiempo Real \\ \hline
Responsable & Gabriel Murillo -- Equipo QA \\ \hline
Fecha de inicio & 15/01/2026 \\ \hline
Fecha de fin & 30/01/2026 \\ \hline
Objetivo & Validar precisión, latencia y robustez gRPC \\ \hline
\end{tabular}
\end{table}

\subsection{2. Alcance}
\begin{table}[h!]
\centering
\begin{tabular}{|l|c|c|}
\hline
\rowcolor{headergrey} \textbf{Tipo de prueba} & \textbf{Incluido} & \textbf{Excluido} \\ \hline
Funcionalidad & \cellcolor{statusgreen} $\checkmark$ & $\times$ \\ \hline
Usabilidad & \cellcolor{statusgreen} $\checkmark$ & $\times$ \\ \hline
Seguridad & \cellcolor{statusgreen} $\checkmark$ & $\times$ \\ \hline
Rendimiento & \cellcolor{statusgreen} $\checkmark$ & $\times$ \\ \hline
\end{tabular}
\end{table}

\subsection{3. Recursos}
\begin{table}[h!]
\centering
\begin{tabular}{|l|c|l|}
\hline
\rowcolor{headergrey} \textbf{Recurso} & \textbf{Cantidad} & \textbf{Observaciones} \\ \hline
Testers & 2 & Evaluación de precisión y latencia \\ \hline
Hardware & 2 & Micrófonos, GPU NVIDIA para Riva \\ \hline
Herramientas & -- & Riva SDK, gRPC, Python, PortAudio \\ \hline
\end{tabular}
\end{table}

\subsection{4. Cronograma}
\begin{table}[h!]
\centering
\begin{tabular}{|c|l|l|}
\hline
\rowcolor{headergrey} \textbf{Semana} & \textbf{Actividad} & \textbf{Responsable} \\ \hline
1 & Diseño y revisión de casos de prueba & QA Lead \\ \hline
2 & Ejecución de pruebas funcionales y gRPC & Testers \\ \hline
3 & Seguimiento de bugs y análisis de causa & QA Lead \\ \hline
4 & Entrega de Informe Final y Cierre & QA Lead \\ \hline
\end{tabular}
\end{table}

\subsection{5. Criterios de Entrada / Salida}
\begin{table}[h!]
\centering
\begin{tabular}{|l|l|}
\hline
\rowcolor{headergrey} \textbf{Criterio} & \textbf{Condición de Aceptación} \\ \hline
Entrada & Código desplegado, NVIDIA Riva activado \\ \hline
Salida & Latencia $\leq$ 500ms, estabilidad gRPC \\ \hline
\end{tabular}
\end{table}

\subsection{6. Ejemplo de Caso de Prueba}
\begin{table}[h!]
\centering
\small
\renewcommand{\arraystretch}{1.3}
\begin{tabular}{|l|l|l|l|l|c|}
\hline
\rowcolor{headergrey} \textbf{ID} & \textbf{Descripción} & \textbf{Pasos} & \textbf{Esperado} & \textbf{Real} & \textbf{Estado} \\ \hline
TC-01 & Conexión gRPC & Iniciar & Canal abierto & Timeout & \cellcolor{statusred} Fallo \\ \hline
TC-02 & Transcripción & Hablar & Texto OK & Texto OK & \cellcolor{statusgreen} Pasó \\ \hline
\end{tabular}
\end{table}

\subsection{7. Métricas de Seguimiento}
\begin{table}[h!]
\centering
\begin{tabular}{|l|c|r|}
\hline
\rowcolor{headergrey} \textbf{Métrica} & \textbf{Valor} & \textbf{Porcentaje} \\ \hline
Casos planificados & 120 & 100\% \\ \hline
Casos ejecutados & 115 & 96\% \\ \hline
Casos pasados & 110 & 91\% \\ \hline
Casos fallidos & 5 & 4\% \\ \hline
\end{tabular}
\end{table}

\subsection{8. Informe Final}
\begin{table}[h!]
\centering
\begin{tabular}{|l|c|}
\hline
\rowcolor{headergrey} \textbf{Elemento} & \textbf{Resultado Final} \\ \hline
Defectos críticos & 0 \\ \hline
Defectos mayores & 2 \\ \hline
Defectos menores & 3 \\ \hline
\rowcolor{statusgreen} \textbf{Conclusión} & \textbf{LISTO PARA PRODUCCIÓN} \\ \hline
\end{tabular}
\end{table}

\section{Actividad para Estudiantes}
Dividir la clase en equipos y entregar esta plantilla. Cada equipo debe completarla simulando un proyecto de IA o procesamiento de señales y presentar el informe final.

\section{Recomendaciones Finales}
\begin{itemize}
  \item Para ASR: mantener lista de frases/hallucinations actualizada.
  \item Para QA: priorizar bloqueantes de red gRPC.
  \item Informes ejecutivos: usar visualizaciones de latencia y WER.
\end{itemize}

\section*{Conclusión}
La gestión de pruebas cubre planificación, ejecución, análisis y comunicación. Este documento y la plantilla asociada están diseñados para uso didáctico en laboratorio y como guía práctica para equipos QA.

\vfill
\noindent Documento generado: 15/01/2026 -- Gabriel Murillo

\end{document}
